%=====================================
%   truncated-ss.tex
%   Microlocal Morse Theory and Truncated Microsupport 
%   toshi2019
%=====================================

% -----------------------
% preamble
% -----------------------
% ここから本文 (\begin{document}) までの
% ソースコードに変更を加えた場合は
% 編集者まで連絡してください. 
% Don't change preamble code yourself. 
% If you add something
% (usepackage, newtheorem, newcommand, renewcommand),
% please tell it 
% to the editor of institutional paper of RUMS.

% ------------------------
% documentclass
% ------------------------
\documentclass[leqno, 10pt, a4paper, dvipdfmx]{article}

% ------------------------
% usepackage
% ------------------------
\usepackage{algorithm}
\usepackage{algorithmic}
\usepackage{amscd}
\usepackage{amsfonts}
\usepackage{amsmath}
\usepackage[psamsfonts]{amssymb}
\usepackage{amsthm}
\usepackage{ascmac}
\usepackage{bm}
\usepackage{color}
\usepackage{enumerate}
\usepackage{fancybox}
\usepackage[stable]{footmisc}
\usepackage{graphicx}
\usepackage{listings}
\usepackage{mathrsfs}
\usepackage{mathtools}
\usepackage{otf}
\usepackage{pifont}
\usepackage{proof}
\usepackage{subfigure}
\usepackage{tikz}
\usepackage{verbatim}
\usepackage[all]{xy}

\usetikzlibrary{cd}
\usetikzlibrary{arrows.meta}


% ================================
% パッケージを追加する場合のスペース 
\usepackage[dvipdfmx]{hyperref}
\usepackage{xcolor}
\definecolor{darkgreen}{rgb}{0,0.45,0} 
\definecolor{darkred}{rgb}{0.75,0,0}
\definecolor{darkblue}{rgb}{0,0,0.6} 
\hypersetup{
    colorlinks=true,
    citecolor=darkgreen,
    linkcolor=darkred,
    urlcolor=darkblue,
}
\usepackage{pxjahyper}

%\usepackage[mathscr]{euscript} % mathscr のフォントを変更

%\usepackage{mmasym}

%=================================


% --------------------------
% theoremstyle
% --------------------------
\theoremstyle{definition}

% --------------------------
% newtheoem
% --------------------------

% 日本語で定理, 命題, 証明などを番号付きで用いるためのコマンドです. 
% If you want to use theorem environment in Japanece, 
% you can use these code. 
% Attention!
% All theorem enivironment numbers depend on 
% only section numbers.
\newtheorem{Axiom}{公理}[section]
\newtheorem{Definition}[Axiom]{定義}
\newtheorem{Theorem}[Axiom]{定理}
\newtheorem{Proposition}[Axiom]{命題}
\newtheorem{Lemma}[Axiom]{補題}
\newtheorem{Corollary}[Axiom]{系}
\newtheorem{Example}[Axiom]{例}
\newtheorem{Claim}[Axiom]{主張}
\newtheorem{Property}[Axiom]{性質}
\newtheorem{Attention}[Axiom]{注意}
\newtheorem{Question}[Axiom]{問}
\newtheorem{Problem}[Axiom]{問題}
\newtheorem{Consideration}[Axiom]{考察}
\newtheorem{Alert}[Axiom]{警告}
\newtheorem{Fact}[Axiom]{事実}


% 日本語で定理, 命題, 証明などを番号なしで用いるためのコマンドです. 
% If you want to use theorem environment with no number in Japanese, You can use these code.
\newtheorem*{Axiom*}{公理}
\newtheorem*{Definition*}{定義}
\newtheorem*{Theorem*}{定理}
\newtheorem*{Proposition*}{命題}
\newtheorem*{Lemma*}{補題}
\newtheorem*{Example*}{例}
\newtheorem*{Corollary*}{系}
\newtheorem*{Claim*}{主張}
\newtheorem*{Property*}{性質}
\newtheorem*{Attention*}{注意}
\newtheorem*{Question*}{問}
\newtheorem*{Problem*}{問題}
\newtheorem*{Consideration*}{考察}
\newtheorem*{Alert*}{警告}
\newtheorem{Fact*}{事実}


% 英語で定理, 命題, 証明などを番号付きで用いるためのコマンドです. 
%　If you want to use theorem environment in English, You can use these code.
%all theorem enivironment number depend on only section number.
\newtheorem{Axiom+}{Axiom}[section]
\newtheorem{Definition+}[Axiom+]{Definition}
\newtheorem{Theorem+}[Axiom+]{Theorem}
\newtheorem{Proposition+}[Axiom+]{Proposition}
\newtheorem{Lemma+}[Axiom+]{Lemma}
\newtheorem{Example+}[Axiom+]{Example}
\newtheorem{Corollary+}[Axiom+]{Corollary}
\newtheorem{Claim+}[Axiom+]{Claim}
\newtheorem{Property+}[Axiom+]{Property}
\newtheorem{Attention+}[Axiom+]{Attention}
\newtheorem{Question+}[Axiom+]{Question}
\newtheorem{Problem+}[Axiom+]{Problem}
\newtheorem{Consideration+}[Axiom+]{Consideration}
\newtheorem{Alert+}{Alert}
\newtheorem{Fact+}[Axiom+]{Fact}
\newtheorem{Remark+}[Axiom+]{Remark}

% ----------------------------
% commmand
% ----------------------------
% 執筆に便利なコマンド集です. 
% コマンドを追加する場合は下のスペースへ. 

% 集合の記号 (黒板文字)
\newcommand{\NN}{\mathbb{N}}
\newcommand{\ZZ}{\mathbb{Z}}
\newcommand{\QQ}{\mathbb{Q}}
\newcommand{\RR}{\mathbb{R}}
\newcommand{\CC}{\mathbb{C}}
\newcommand{\PP}{\mathbb{P}}
\newcommand{\KK}{\mathbb{K}}


% 集合の記号 (太文字)
\newcommand{\nn}{\mathbf{N}}
\newcommand{\zz}{\mathbf{Z}}
\newcommand{\qq}{\mathbf{Q}}
\newcommand{\rr}{\mathbf{R}}
\newcommand{\cc}{\mathbf{C}}
\newcommand{\pp}{\mathbf{P}}
\newcommand{\kk}{\mathbf{k}}

% 特殊な写像の記号
\newcommand{\ev}{\mathop{\mathrm{ev}}\nolimits} % 値写像
\newcommand{\pr}{\mathop{\mathrm{pr}}\nolimits} % 射影

% スクリプト体にするコマンド
%   例えば　{\mcal C} のように用いる
\newcommand{\mcal}{\mathcal}

% 花文字にするコマンド 
%   例えば　{\h C} のように用いる
\newcommand{\h}{\mathscr}

% ヒルベルト空間などの記号
\newcommand{\F}{\mcal{F}}
\newcommand{\X}{\mcal{X}}
\newcommand{\Y}{\mcal{Y}}
\newcommand{\Hil}{\mcal{H}}
\newcommand{\RKHS}{\Hil_{k}}
\newcommand{\Loss}{\mcal{L}_{D}}
\newcommand{\MLsp}{(\X, \Y, D, \Hil, \Loss)}

% 偏微分作用素の記号
\newcommand{\p}{\partial}

% 角カッコの記号 (内積は下にマクロがあります)
\newcommand{\lan}{\langle}
\newcommand{\ran}{\rangle}



% 圏の記号など
\newcommand{\Set}{{\bf Set}}
\newcommand{\Vect}{{\bf Vect}}
\newcommand{\FDVect}{{\bf FDVect}}
%\newcommand{\Ring}{{\bf Ring}}
\newcommand{\Ab}{{\bf Ab}}
\newcommand{\Mod}{\mathop{\mathrm{Mod}}\nolimits}
\newcommand{\Modf}{\mathop{\mathrm{Mod}^\mathrm{f}}\nolimits}
\newcommand{\CGA}{{\bf CGA}}
\newcommand{\GVect}{{\bf GVect}}
\newcommand{\Lie}{{\bf Lie}}
\newcommand{\dLie}{{\bf Liec}}



% 射の集合など
\newcommand{\Map}{\mathop{\mathrm{Map}}\nolimits} % 写像の集合
\newcommand{\Hom}{\mathop{\mathrm{Hom}}\nolimits} % 射集合
\newcommand{\End}{\mathop{\mathrm{End}}\nolimits} % 自己準同型の集合
\newcommand{\Aut}{\mathop{\mathrm{Aut}}\nolimits} % 自己同型の集合
\newcommand{\Mor}{\mathop{\mathrm{Mor}}\nolimits} % 射集合
\newcommand{\Ker}{\mathop{\mathrm{Ker}}\nolimits} % 核
\newcommand{\Img}{\mathop{\mathrm{Im}}\nolimits} % 像
\newcommand{\Cok}{\mathop{\mathrm{Coker}}\nolimits} % 余核
\newcommand{\Cim}{\mathop{\mathrm{Coim}}\nolimits} % 余像

% その他便利なコマンド
\newcommand{\dip}{\displaystyle} % 本文中で数式モード
\newcommand{\e}{\varepsilon} % イプシロン
\newcommand{\dl}{\delta} % デルタ
\newcommand{\pphi}{\varphi} % ファイ
\newcommand{\ti}{\tilde} % チルダ
\newcommand{\pal}{\parallel} % 平行
\newcommand{\op}{{\rm op}} % 双対を取る記号
\newcommand{\lcm}{\mathop{\mathrm{lcm}}\nolimits} % 最小公倍数の記号
\newcommand{\Probsp}{(\Omega, \F, \P)} 
\newcommand{\argmax}{\mathop{\rm arg~max}\limits}
\newcommand{\argmin}{\mathop{\rm arg~min}\limits}





% ================================
% コマンドを追加する場合のスペース 
%\newcommand{\OO}{\mcal{O}}



\makeatletter
\renewenvironment{proof}[1][\proofname]{\par
  \pushQED{\qed}%
  \normalfont \topsep6\p@\@plus6\p@\relax
  \trivlist
  \item[\hskip\labelsep
%        \itshape
         \bfseries
%    #1\@addpunct{.}]\ignorespaces
    {#1}]\ignorespaces
}{%
  \popQED\endtrivlist\@endpefalse
}
\makeatother

\renewcommand{\proofname}{Proof.}



%\renewcommand\proofname{\bf 証明} % 証明
\numberwithin{equation}{section}
\newcommand{\cTop}{\textsf{Top}}
%\newcommand{\cOpen}{\textsf{Open}}
\newcommand{\Op}{\mathop{\textsf{Op}}\nolimits}
\newcommand{\Ob}{\mathop{\textrm{Ob}}\nolimits}
\newcommand{\id}{\mathop{\mathrm{id}}\nolimits}
\newcommand{\pt}{\mathop{\mathrm{pt}}\nolimits}
\newcommand{\res}{\mathop{\rho}\nolimits}
\newcommand{\A}{\mcal{A}}
\newcommand{\B}{\mcal{B}}
\newcommand{\C}{\mcal{C}}
\newcommand{\D}{\mcal{D}}
\newcommand{\E}{\mcal{E}}
\newcommand{\G}{\mcal{G}}
%\newcommand{\H}{\mcal{H}}
\newcommand{\I}{\mcal{I}}
\newcommand{\J}{\mcal{J}}
\newcommand{\OO}{\mcal{O}}
\newcommand{\Ring}{\mathop{\textsf{Ring}}\nolimits}
\newcommand{\cAb}{\mathop{\textsf{Ab}}\nolimits}
%\newcommand{\Ker}{\mathop{\mathrm{Ker}}\nolimits}
\newcommand{\im}{\mathop{\mathrm{Im}}\nolimits}
\newcommand{\Coker}{\mathop{\mathrm{Coker}}\nolimits}
\newcommand{\Coim}{\mathop{\mathrm{Coim}}\nolimits}
\newcommand{\rank}{\mathop{\mathrm{rank}}\nolimits}
\newcommand{\Ht}{\mathop{\mathrm{Ht}}\nolimits}
\newcommand{\supp}{\mathop{\mathrm{supp}}\nolimits}
\newcommand{\colim}{\mathop{\mathrm{colim}}}
\newcommand{\Tor}{\mathop{\mathrm{Tor}}\nolimits}

\newcommand{\cat}{\mathscr{C}}

%筆記体
\newcommand{\cA}{\mcal{A}}
\newcommand{\cB}{\mcal{B}}
\newcommand{\cC}{\mcal{C}}
\newcommand{\cD}{\mcal{D}}
\newcommand{\cE}{\mcal{E}}
\newcommand{\cF}{\mcal{F}}
\newcommand{\cG}{\mcal{G}}
\newcommand{\cH}{\mcal{H}}
\newcommand{\cI}{\mcal{I}}
\newcommand{\cJ}{\mcal{J}}
\newcommand{\cK}{\mcal{K}}
\newcommand{\cL}{\mcal{L}}
\newcommand{\cM}{\mcal{M}}
\newcommand{\cN}{\mcal{N}}
\newcommand{\cO}{\mcal{O}}
\newcommand{\cP}{\mcal{P}}
\newcommand{\cQ}{\mcal{Q}}
\newcommand{\cR}{\mcal{R}}
\newcommand{\cS}{\mcal{S}}
\newcommand{\cT}{\mcal{T}}
\newcommand{\cU}{\mcal{U}}
\newcommand{\cV}{\mcal{V}}
\newcommand{\cW}{\mcal{W}}
\newcommand{\cX}{\mcal{X}}
\newcommand{\cY}{\mcal{Y}}
\newcommand{\cZ}{\mcal{Z}}


\newcommand{\scA}{\mathscr{A}}
\newcommand{\scB}{\mathscr{B}}
\newcommand{\scC}{\mathscr{C}}
\newcommand{\scD}{\mathscr{D}}
\newcommand{\scE}{\mathscr{E}}
\newcommand{\scF}{\mathscr{F}}
\newcommand{\scN}{\mathscr{N}}
\newcommand{\scO}{\mathscr{O}}
\newcommand{\scV}{\mathscr{V}}
\newcommand{\scU}{\mathscr{U}}


\newcommand{\ibA}{\mathop{\text{\textit{\textbf{A}}}}}
\newcommand{\ibB}{\mathop{\text{\textit{\textbf{B}}}}}
\newcommand{\ibC}{\mathop{\text{\textit{\textbf{C}}}}}
\newcommand{\ibD}{\mathop{\text{\textit{\textbf{D}}}}}
\newcommand{\ibE}{\mathop{\text{\textit{\textbf{E}}}}}
\newcommand{\ibF}{\mathop{\text{\textit{\textbf{F}}}}}
\newcommand{\ibG}{\mathop{\text{\textit{\textbf{G}}}}}
\newcommand{\ibH}{\mathop{\text{\textit{\textbf{H}}}}}
\newcommand{\ibI}{\mathop{\text{\textit{\textbf{I}}}}}
\newcommand{\ibJ}{\mathop{\text{\textit{\textbf{J}}}}}
\newcommand{\ibK}{\mathop{\text{\textit{\textbf{K}}}}}
\newcommand{\ibL}{\mathop{\text{\textit{\textbf{L}}}}}
\newcommand{\ibM}{\mathop{\text{\textit{\textbf{M}}}}}
\newcommand{\ibN}{\mathop{\text{\textit{\textbf{N}}}}}
\newcommand{\ibO}{\mathop{\text{\textit{\textbf{O}}}}}
\newcommand{\ibP}{\mathop{\text{\textit{\textbf{P}}}}}
\newcommand{\ibQ}{\mathop{\text{\textit{\textbf{Q}}}}}
\newcommand{\ibR}{\mathop{\text{\textit{\textbf{R}}}}}
\newcommand{\ibS}{\mathop{\text{\textit{\textbf{S}}}}}
\newcommand{\ibT}{\mathop{\text{\textit{\textbf{T}}}}}
\newcommand{\ibU}{\mathop{\text{\textit{\textbf{U}}}}}
\newcommand{\ibV}{\mathop{\text{\textit{\textbf{V}}}}}
\newcommand{\ibW}{\mathop{\text{\textit{\textbf{W}}}}}
\newcommand{\ibX}{\mathop{\text{\textit{\textbf{X}}}}}
\newcommand{\ibY}{\mathop{\text{\textit{\textbf{Y}}}}}
\newcommand{\ibZ}{\mathop{\text{\textit{\textbf{Z}}}}}

\newcommand{\ibx}{\mathop{\text{\textit{\textbf{x}}}}}

%\newcommand{\Comp}{\mathop{\mathrm{C}}\nolimits}
%\newcommand{\Komp}{\mathop{\mathrm{K}}\nolimits}
%\newcommand{\Domp}{\mathop{\mathsf{D}}\nolimits}%複体のホモトピー圏
%\newcommand{\Comp}{\mathrm{C}}
%\newcommand{\Komp}{\mathrm{K}}
%\newcommand{\Domp}{\mathsf{D}}%複体のホモトピー圏

\newcommand{\Comp}{\mathop{\mathrm{C}}\nolimits}
\newcommand{\Komp}{\mathop{\mathsf{K}}\nolimits}
\newcommand{\Domp}{\mathord{\mathsf{D}}}%\nolimits}
\newcommand{\Kompl}{\mathop{\mathsf{K}^\mathrm{+}}\nolimits}
\newcommand{\Kompu}{\mathop{\mathsf{K}^\mathrm{-}}\nolimits}
\newcommand{\Kompb}{\mathop{\mathsf{K}^\mathrm{b}}\nolimits}
\newcommand{\Dompl}{\mathop{\mathsf{D}^\mathrm{+}}\nolimits}
\newcommand{\Dompu}{\mathop{\mathsf{D}^\mathrm{-}}\nolimits}
\newcommand{\Dompb}{\mathop{\mathsf{D}^\mathrm{b}}\nolimits}
\newcommand{\Dompbf}{\mathop{\mathsf{D}_\mathrm{f}^\mathrm{b}}\nolimits}
\newcommand{\Domplb}{\mathop{\mathsf{D}^\mathrm{lb}}\nolimits}




\newcommand{\CCat}{\Comp(\cat)}
\newcommand{\KCat}{\Komp(\cat)}
\newcommand{\DCat}{\Domp(\cat)}%圏Cの複体のホモトピー圏
\newcommand{\HOM}{\mathop{\mathscr{H}\hspace{-2pt}om}\nolimits}%内部Hom
\newcommand{\RHOM}{\mathop{\mathrm{R}\hspace{-1.5pt}\HOM}\nolimits}

%\newcommand{\muS}{\mathop{\mathrm{SS}}\nolimits}
\newcommand{\muS}[1][]{\mathop{\mathrm{SS}}\nolimits_{#1}}
\newcommand{\RG}{\mathop{\mathrm{R}\hspace{-0pt}\Gamma}\nolimits}
\newcommand{\RHom}{\mathop{\mathrm{R}\hspace{-1.5pt}\Hom}\nolimits}
%\newcommand{\Rder}{\mathrm{R}}
\newcommand{\Rder}[1][]{\mathop{\mathrm{R}^{#1}\hspace{-2pt}}\nolimits}

\newcommand{\simar}{\mathrel{\overset{\sim}{\rightarrow}}}%同型右矢印
\newcommand{\simarr}{\mathrel{\overset{\sim}{\longrightarrow}}}%同型右矢印
\newcommand{\simra}{\mathrel{\overset{\sim}{\leftarrow}}}%同型左矢印
\newcommand{\simrra}{\mathrel{\overset{\sim}{\longleftarrow}}}%同型左矢印

\newcommand{\hocolim}{{\mathrm{hocolim}}}
\newcommand{\indlim}[1][]{\mathop{\varinjlim}\limits_{#1}}
\newcommand{\sindlim}[1][]{\smash{\mathop{\varinjlim}\limits_{#1}}\,}
\newcommand{\Pro}{\mathrm{Pro}}
\newcommand{\Ind}{\mathrm{Ind}}
\newcommand{\prolim}[1][]{\mathop{\varprojlim}\limits_{#1}}
\newcommand{\sprolim}[1][]{\smash{\mathop{\varprojlim}\limits_{#1}}\,}

\newcommand{\Sh}{\mathrm{Sh}}
\newcommand{\PSh}{\mathrm{PSh}}

\newcommand{\rmD}{\mathrm{D}}

\newcommand{\Lloc}[1][]{\mathord{\mathcal{L}^1_{\mathrm{loc},{#1}}}}
\newcommand{\ori}{\mathord{\mathrm{or}}}
\newcommand{\Db}{\mathord{\mathscr{D}b}}

\newcommand{\codim}{\mathop{\mathrm{codim}}\nolimits}



\newcommand{\gld}{\mathop{\mathrm{gld}}\nolimits}
\newcommand{\wgld}{\mathop{\mathrm{wgld}}\nolimits}


\newcommand{\tens}[1][]{\mathbin{\otimes_{\raise1.5ex\hbox to-.1em{}{#1}}}}
\newcommand{\etens}{\mathbin{\boxtimes}}
\newcommand{\ltens}[1][]{\mathbin{\overset{\mathrm{L}}\tens}_{#1}}
\newcommand{\mtens}[1][]{\mathbin{\overset{\mathrm{\mu}}\tens}_{#1}}
\newcommand{\lltens}[1][]{{\mathop{\tens}\limits^{\mathrm{L}}_{#1}}}
\newcommand{\letens}{\overset{\mathrm{L}}{\etens}}
\newcommand{\detens}{\underline{\etens}}
\newcommand{\ldetens}{\overset{\mathrm{L}}{\underline{\etens}}}
\newcommand{\dtens}[1][]{{\overset{\mathrm{L}}{\underline{\otimes}}}_{#1}}

\newcommand{\blk}{\mathord{\ \cdot\ }}
\newcommand{\mres}[2][]{{\left.{#1}\right\rvert}_{#2}}


%\newcommand{\hocolim}{{\mathrm{hocolim}}}
%\newcommand{\indlim}[1][]{\mathop{\varinjlim}\limits_{#1}}
%\newcommand{\sindlim}[1][]{\smash{\mathop{\varinjlim}\limits_{#1}}\,}
%\newcommand{\Pro}{\mathrm{Pro}}
%\newcommand{\Ind}{\mathrm{Ind}}
%\newcommand{\prolim}[1][]{\mathop{\varprojlim}\limits_{#1}}
%\newcommand{\sprolim}[1][]{\smash{\mathop{\varprojlim}\limits_{#1}}\,}
\newcommand{\proolim}[1][]{\mathop{\text{\rm``{$\varprojlim$}''}}\limits_{#1}}
\newcommand{\sproolim}[1][]{\smash{\mathop{\rm``{\varprojlim}''}\limits_{#1}}}
\newcommand{\inddlim}[1][]{\mathop{\text{\rm``{$\varinjlim$}''}}\limits_{#1}}
\newcommand{\sinddlim}[1][]{\smash{\mathop{\text{\rm``{$\varinjlim$}''}}\limits_{#1}}\,}
\newcommand{\ooplus}{\mathop{\text{\rm``{$\oplus$}''}}\limits}
\newcommand{\bbigsqcup}{\mathop{``\bigsqcup''}\limits}
\newcommand{\bsqcup}{\mathop{``\sqcup''}\limits}
\newcommand{\dsum}[1][]{\mathbin{\oplus_{#1}}}

\newcommand{\Fct}{\mathop{\mathsf{Fct}}\nolimits}

\newcommand{\pb}[1][]{\mathbin{\mathop{\times}\limits_{#1}}}
\newcommand{\dpb}[1][]{\mathbin{\mathop{\times}\limits_{#1}}}
\newcommand{\tpb}[1][]{\mathbin{\mathop{\times}\nolimits_{#1}}}





%================================================
% 自前の定理環境
%   https://mathlandscape.com/latex-amsthm/
% を参考にした
\newtheoremstyle{mystyle}%   % スタイル名
    {5pt}%                   % 上部スペース
    {5pt}%                   % 下部スペース
    {}%              % 本文フォント
    {}%                  % 1行目のインデント量
    {\bfseries}%                      % 見出しフォント
    {.}%                     % 見出し後の句読点
    {12pt}%                     % 見出し後のスペース
    {\thmname{#1}\thmnumber{ #2}\thmnote{{\hspace{2pt}\normalfont (#3)}}}% % 見出しの書式

\theoremstyle{mystyle}
\newtheorem{AXM}{公理}[section]
\newtheorem{DFN}[AXM]{定義}
\newtheorem{THM}[AXM]{定理}
\newtheorem*{THM*}{定理}
\newtheorem{PRP}[AXM]{命題}
\newtheorem{LMM}[AXM]{補題}
\newtheorem{CRL}[AXM]{系}
\newtheorem{EG}[AXM]{例}
\newtheorem*{EG*}{例}
\newtheorem{RMK}[AXM]{注意}
\newtheorem{CNV}[AXM]{約束}
\newtheorem{CMT}[AXM]{コメント}
\newtheorem*{CMT*}{コメント}
\newtheorem{NTN}[AXM]{記号}

% 定理環境ここまで
%====================================================

\usepackage{framed}
\definecolor{lightgray}{rgb}{0.75,0.75,0.75}
\renewenvironment{leftbar}{%
  \def\FrameCommand{\textcolor{lightgray}{\vrule width 4pt} \hspace{10pt}}% 
  \MakeFramed {\advance\hsize-\width \FrameRestore}}%
{\endMakeFramed}
\newenvironment{redleftbar}{%
  \def\FrameCommand{\textcolor{lightgray}{\vrule width 1pt} \hspace{10pt}}% 
  \MakeFramed {\advance\hsize-\width \FrameRestore}}%
 {\endMakeFramed}



\renewcommand{\Re}{\mathop{\mathrm{Re}}\nolimits}

% =================================





% ---------------------------
% new definition macro
% ---------------------------
% 便利なマクロ集です

% 内積のマクロ
%   例えば \inner<\pphi | \psi> のように用いる
\def\inner<#1>{\langle #1 \rangle}

% ================================
% マクロを追加する場合のスペース 

%=================================



% 位相空間
\newcommand{\Int}[1][]{\mathop{\mathrm{Int}}\nolimits_{#1}}
\newcommand{\Cl}[1][]{\mathop{\mathrm{Cl}}\nolimits_{#1}}
\newcommand{\Bd}[1][]{\mathop{\mathrm{Bd}}\nolimits_{#1}}
\newcommand{\bd}{\mathord{\partial}}
\newcommand{\Fr}[1][]{\mathop{\mathrm{Fr}}\nolimits_{#1}}


\newcommand{\mtan}[1][]{T\hspace{-1.75pt}{#1}}
\newcommand{\mtp}[1][]{{}^{t\!}{#1}} % transpose of the morphism
%\newcommand{\mpb}[1][]{{}^{t\!}{#1}'} % pull-back of the morphism
\newcommand{\mpb}[1][]{{#1}_{\pi}} % pull-back of the morphism
\newcommand{\mdiff}[1][]{{#1}_{d}} % pull-back of the morphism
\newcommand{\mptan}[1][]{{#1}_{\tau}} % pull-back of the morphism


%\newcommand{\pb}[1][]{\mathbin{\mathop{\times}\limits_{#1}}}

\newcommand{\cotb}[1][]{T^{\ast}{#1}}
\newcommand{\conm}[2][]{T_{#1}^{\hspace{0.7pt}\ast}{#2}}
\newcommand{\norb}[2][]{T_{#1}{#2}}

\newcommand{\cotz}[1][]{T^{\ast}_{#1}{#1}}

\newcommand{\cotn}[1][]{\dot{T}^{\ast}{#1}}


\newcommand{\muhom}{\mu hom}
\newcommand{\muHom}[1][]{\mathrm{Hom}^\mu_{\raise1.5ex\hbox to.1em{}#1}}

\newcommand{\mucat}[2][]{\mathsf{D}^{\mathrm{b}}_{#1}(#2)}
%\newcommand{\mucat}[2][]{\mathsf{D}^{\mathrm{b}}_{#1}(#2)}

\newcommand{\conv}[1][]{\mathbin{\mathop{\circ}\limits_{#1}}}


\newcommand{\lopen}{\mathopen{]}}
\newcommand{\lclosed}{\mathopen{[}}
\newcommand{\ropen}{\mathclose{[}}
\newcommand{\rclosed}{\mathopen{]}}

\newcommand{\sm}{\raisebox{2.33pt}{~\rule{6.4pt}{1.3pt}~}}
\newcommand{\smi}{\mathbin{\raisebox{2.33pt}{\rule{6.4pt}{1.3pt}}}}

\newcommand{\smid}{\mathrel{;}}


% ----------------------------
% documenet 
% ----------------------------
% 以下, 本文の執筆スペースです. 
% Your main code must be written between 
% begin document and end document.
% ---------------------------

\title{Microlocal Morse Theory and Truncated Microsupport}
\author{Toshihiro Oshiba}
\date{\today\footnote{2025/03/09 first draft}}
\begin{document}
\maketitle
%\tableofcontents
%\tableofcontents
\section*{Preface}
In these notes, 
We study microlocal Morse theory, which is presented, 
for instance, in \cite{GKS12}. 
However, we replace the notion of microsupports with
 that of truncatd microsupprts, introduced by 
Kashiwara, Monteiro-Fernandes, and Schapira in \cite{KMS03}.

\section{Notations}

We consider a commutative unital ring \(\kk\) 
of finite global dimension (e.g. \(\zz\) or a field). 
Let \(X\) be a real analytic manifold. 
We denote the cotangent bundle of \(X\) by \(
  \pi\colon \cotb[X]\to X
\). 
We denote by \(\Dompb(\kk_X)\) the bounded derived category of 
sheaves of \(\kk\)-modules on \(X\). 

\section{Truncated Microsupport}

\begin{Definition+}
  Let \(F\in\Dompb(\kk_X)\). 
  The closed conic subset \(\muS[k](F)\) of \(\cotb[X]\) 
  is defined by: \(p\notin \muS[k](F)\) 
  if and only if 
  there exists an open conic neighborhood \(U\) of \(p\) 
  such that for any \(x\in \pi(U)\) and 
  for any real-valued \(C^{1}\)-function \(\varphi\) 
  defined on a neighborhood of \(x\) such that 
  \(\varphi(x)=0\), \(d\varphi(x)\in U\), 
  one has \begin{equation}
    H^{j}_{\{\varphi\geqq0\}}(F)_x=0
    \quad \text{for any \(j\leqq k\).}
  \end{equation}
\end{Definition+}

\subsection{Functorial Properties} 
\begin{Proposition+}
  Let \(f\colon X\to Y\) be a morphism of real analytic manifolds and let \(F\in\Dompb(\kk_X)\) such that \(f\) is proper on \(\supp(F)\). 
  Then for any \(k\in\zz\), 
  \begin{equation}\label{eq:t-push}
    \muS[k](\Rder{f_{\ast}}F)\subset
    \mpb[f]\mdiff[f]^{-1}(\muS[k](F)).
  \end{equation}
  The equality holds in case \(f\) is a closed embedding.
\end{Proposition+}

\section{Lemmas from {\cite{KS90}}}

Consider a projective system of abelian groups \(
  \left(M_n,\rho_{n,p}\right)_n
\) indexed by \(\nn\). 
One says that this system satisfies 
the Mittag-Leffler condition (ML for short) 
if for any \(n\in\nn\), the decreasing sequence \(
  (\rho_{n,p}(M_p))_p
\) of subgroups of \(M_n\) is stationary.

For a projective system of abelian groups \((M_{n})_n\), 
we set \[
  M_{\infty}\coloneqq \prolim[n]M_n.
\]

Consider a projective system of complexes
\begin{equation}
  E_n^{\bullet}\colon 
  \cdots\to M_{n}^{j-1}\to M_{n}^{j}\to M_{n}^{j+1}\to\cdots,
\end{equation}
and projective limit
\begin{equation}
  E_{\infty}^{\bullet}\colon 
  \cdots
  \to M_{\infty}^{j-1}
  \to M_{\infty}^{j}
  \to M_{\infty}^{j+1}
  \to\cdots.
\end{equation}
Denote by 
\begin{equation}
  \varPhi_{k}\colon H^{k}(E_{\infty}^{\bullet})\to \prolim[n]H^{k}(E_{n}^{\bullet})
\end{equation}
the natural morphism.

\begin{Proposition+}[{\cite[Proposition 1.12.4]{KS90}}]\label{prop:ks1.12.4}
  Assume that for each \(j\in\zz\), 
  the system \(\left(M_{n}^{j}\right)_{n}\) satisfies 
  the ML condition. 
  Then
  \begin{enumerate}[(a)]
    \item for each \(k\in \zz\), the map \(\varPhi_k\) is surjective,
    \item if moreover, for a given \(i\) the system \(\left(H^{i-1}(E_{n}^{\bullet})\right)_{n}\) satisfies the ML condition, then \(\varPhi_i\) is bijective.
  \end{enumerate} 
\end{Proposition+}

The following is called Kashiwara's lemma.

\begin{Proposition+}[{\cite[Proposition 1.12.6]{KS90}}]\label{prop:ks1.12.6}
  Let \(\left(X_s,\rho_{s,t}\colon X_t\to X_s\right)_{(s,(s\leqq t))}\) be 
  a projective system of sets indexed by \(\rr\). 
  Assume that for each \(s\in\rr\), the canonical maps\[
    \lambda_s\colon X_s\to\prolim[r<s]X_r
  \] and \[
    \mu_s\colon \indlim[t>s]X_t\to X_s
  \] are both injective (resp. surjective). 
  Then all the maps \(\rho_{s_{0},s_{1}}\) \((s_0\leqq s_1)\) are 
  injective (resp. surjective).
\end{Proposition+}

\subsection*{ML procedure for sheaves}

\begin{Proposition+}[{\cite[Proposition 2.7.1]{KS90}}]\label{prop:ks2.7.1}
  Let \(X\) be a topological space and let \(F\in\Dompl(\zz_X)\). 
  Let \(\left(U_n\right)_{n\in\nn}\) be an increasing sequence 
  of open subsets of \(X\), \(\left(Z_n\right)_{n\in\nn}\) 
  a decreasing sequence of closed subsets of \(X\). 
  Set \(U\coloneqq\bigcup_{n}U_{n}\), 
  \(Z\coloneqq\bigcap_{n}Z_{n}\). 
  \begin{enumerate}[(i)]
    \item For any \(j\), the natural map \(
      \varphi_j\colon H^{j}_{Z}(U;F)
      \to 
      \prolim[n]H^{j}_{Z_{n}}(U_{n};F)
    \) is surjective. 
    \item Assume that for a given \(j\), 
    the projective system \(
      (H^{j-1}_{Z_{n}}(U_{n};F))_{n}
    \) satisfies the ML condition. 
    Then \(\varphi_{j}\) is bijective. 
    \item Let now \(\left(X_{n}\right)_{n\in\nn}\) be 
    an increasing family of subsets of \(X\) satisfying \(
      X=\bigcup_{n}X_n\) and \(X_{n}\subset \Int[](X_{n+1})
    \) for all \(n\). 
    Assume that for a given \(j\), the projective system \(
      \left(H^{j-1}(X_{n};F)\right)_{n}
    \) satisfies the ML condition. 
    Then the natural map \(
      H^{j}(X;F)\to \prolim[n]H^{j}(X_{n};F)
    \) is bijective. 
  \end{enumerate}
\end{Proposition+}

\section{Results}

In this section, 
let \(M\) be a manifold.\footnote{
  Can \(M\) be of class \(C^\infty\)?
}

\begin{Theorem+}\label{thm:t-morse}
  Let \(F\) in \(\Dompb(\kk_M)\), and 
  \(\psi\colon M\to \rr\) be a real-valued function of 
  class \(C^1\), 
  proper on \(\supp(F)\). 
  Let \(a<b\) be real numbers.
  We assume \(d\psi(x)\notin \muS[k](F)\) for \(
    a\leqq \psi(x)<b
  \). For \(t\in \rr\), we set \(
    M_t\coloneqq\psi^{-1}\left(]-\infty,t[\right)
  \). 
  Then, the restriction morphisms\[
    H^{j}(M_{b};F)\to H^{j}(M_{a};F)
  \]
  are isomorphic for each \(j< k\).\footnote{
    For \(j=k\), this morphism might be injective.
  }
\end{Theorem+}

\begin{Lemma+}\label{Lem:1d-t-morse}
  Let \(-\infty<a<b<+\infty\) be real numbers 
  and \(G\in \Dompb(\rr)\). Let \(k\) be an integer.
  If \[
      \left(H^{j}_{\left[t,+\infty\right[}(G)\right)_t
      \cong 0
      \quad
      \text{
          for all \(t\in\left[a,b\right[\) , 
          and all \(j\leqq k\).
      }
  \]
  Then one has \[
      H^{j}(\left]-\infty,b\right[;G)\simarr
      H^{j}(\left]-\infty,a\right[;G)
  \] for all \(j<k\).\footnote{
      For \(k\), this morphism might be injective.
        (not written in this draft.)
  }
\end{Lemma+}

\begin{proof}
  Let \(t\in \lclosed a,b\ropen\). 
  We will denote by \(I_t\) the open interval \(
      \left]-\infty,t\right[
  \) and denote by \(Z_t\) the closed interval \(
      \left]-\infty,t\right] (= \overline{I_t})
  \).
  By applying the functor \((\blk)_{\overline{I_t}}\) 
  to the distinguished triangle\[
      \RG_{\rr-I_t}(G)\to 
      G\to 
      \RG_{I_t}(G)\overset{+1}{\longrightarrow}
      \quad\text{ in \(\Dompb(\kk_\rr)\)},
  \] we have\[
      (\RG_{[t,+\infty[}(G))_{\{t\}}
      \to G_{Z_t}
      \to 
      \RG_{I_t}(G)
      \overset{+1}{\longrightarrow}
      \quad\text{ in \(\Dompb(\kk_\rr)\)}.
  \]
  Aplying the functor \(\RG(\rr;\blk)\) to this, we have
  \begin{equation}\label{eq:morse-d.t.}
      (\RG_{[t,+\infty[}(G))_{t}
      \to \RG(Z_t;G)
      \to 
      \RG(I_t;G)
      \overset{+1}{\longrightarrow}
      \quad\text{ in \(\Dompb(\kk)\)}.
  \end{equation}
  Then consider the following conditions for 
  \(j\in\zz\), \(s\in\left[a,b\right[\), and 
  \(t\in \left]a,b\right]\):
  \begin{align}
      \tag*{\((a)^j_s\)}\label{eq:nonchar-a}
      \indlim[t>s] H^{j}(I_t;G) &\simarr H^{j}(I_s;G),\\
      \tag*{\((b)^j_t\)}\label{eq:nonchar-b}
      \prolim[s<t] H^{j}(I_t;G) &\simrra H^{j}(I_t;G).
  \end{align}
  By hypothesis, 
  it follows from \eqref{eq:morse-d.t.} that 
  \(
    H^{j}(Z_t;G)\simar H^{j}(I_t;G)
  \) for \(t\in \left[a,b\right[\).\footnote{
    Similarly, we have injections \(
      H^{k}(Z_t;G)\hookrightarrow H^{k}(I_t;G)
  \) for \(t\in \left[a,b\right[\). 
  } 
  Since \(
      \indlim[t>s] H^{j}(I_t;G)\cong H^{j}(Z_s;G)
  \), \ref{eq:nonchar-a} holds 
  for any \(j<k\) and \(s\in\left[a,b\right[\).\footnote{
      Similarly, \[
        \indlim[t>s] H^{k}(I_t;G)\hookrightarrow H^{k}(Z_s;G)
    \] gives injections \[
        \indlim[t>s] H^{k}(I_t;G)\hookrightarrow H^{k}(I_s;G)
    \] for all \(s\in \left[a,b\right[\). 
  } Moreover, since we assume \(G\) is a bounded complex, 
  \ref{eq:nonchar-b} holds for \(j\ll0\) 
  and for all \(t\in \left]a,b\right]\). 
  Let us denote such \(j\) by \(j_0\). 
  Let us argue by induction on \(j\). 
  Assume \ref{eq:nonchar-b} holds for all 
  \(t\in\left]a,b\right[\) and all \(j\leqq j_0\). 
  Applying Proposition \ref{prop:ks1.12.6}, 
  we find the isomorphisms \begin{equation}\label{eq:non-char-init}
    H^{j}(I_s;G) \simarr H^{j}(I_t;G)
    \quad
    \text{for all \(j\leqq j_0\) and all \(a<s\leqq t< b\).}
  \end{equation}
  Let \(G \simar G^{\bullet}\) be a flabby resolution. 
  Let \(t\in\left]a,b\right]\). 
  Consider a complex \[
    E^{\bullet}_{n}=\Gamma(I_{t-\frac{1}{n}};G^{\bullet}).
  \]
  Since \(G^{\bullet}\) is flabby, 
  we have \(E_{n}^{j}\cong E_{\infty}^{j}\) and 
  the projective system \((E^{j}_{n})_{n}\) satisfies 
  the ML condition for all \(j\in\zz\). 
  
  By \eqref{eq:non-char-init}, the projective system \(
    (H^{j_0}(E^{\bullet}_{n}))_{n}
  \) satisfies the ML condition (because we have \(
    a<t-\frac{1}{n}<t
  \) for \(n\gg0\)). 
  Applying Proposition \ref{prop:ks1.12.4} to \(
    (E^{\bullet}_{n})_{n}
  \), it follows that \ref{eq:nonchar-b} holds for \(j_{0}+1\) 
  and the induction proceeds. More precisely, we have \begin{align*}
    H^{j_0+1}(I_{t};G^{\bullet})
    \overunderset{\sim}{\varPhi_{j_{0}+1}}{\longrightarrow}&
    \prolim[n]H^{j_0+1}(I_{t-\frac{1}{n}};G^{\bullet})\\
    \cong&
    \prolim[n]H^{j_0+1}(I_{t-\frac{1}{n}};G)\\
    \cong&
    \prolim[s<t]H^{j_0+1}(I_{s};G).
  \end{align*}
  Again by Proposition \ref{prop:ks1.12.6}, 
  we get the isomorphisms\[
    H^{j}(I_{s};G)\simarr H^{j}(I_t;G)
    \quad
    \text{for all \(j<k\) and all \(a<s\leqq t\leqq b\),}
  \]
  and injections \[
    H^{j}(I_{s};G)\hookrightarrow H^{j}(I_t;G)
  \] for \(j=k\) and for all \(a<s\leqq t\leqq b\). 

  Finally, using \eqref{eq:morse-d.t.} and the hypothesis, 
  we have the isomorphisms\[
    H^{j}(I_{a};G)\simrra
    H^{j}(Z_{a};G)\simarr H^{j}(I_s;G)
    \quad
    \text{for all \(j<k\) and all \(a<s\leqq b\).}
  \]
  Similarly, for \(j=k\), we have\[
    H^{k}(I_{b};G) \cong \indlim[t>s] H^{k}(I_t;G)\hookrightarrow H^{k}(I_s;G) \hookrightarrow H^{k}(I_{a};G).
  \]
\end{proof}

\begin{proof}[Proof of Theorem {\ref{thm:t-morse}}.]
  First, we have\[
    H^{j}(M_{t};F)
    \cong
    H^{j}\RG(
      {\left]-\infty,t\right[}\mathrel{;}\Rder{\psi}_{\ast}F
    )
  \]
  for \(t\in I\).
  Since we assume that \(d\psi(x)\notin \muS[k](F)\) 
  for \(a\leqq \psi(x)<b\), and by \eqref{eq:t-push}, the formula \(
    \muS[k](\Rder{\psi}_{\ast}F)
    \subset \mpb[\psi]\mdiff[\psi]^{-1}\muS[k](F)
  \) holds. Hence we have\[
    \left(H^{j}\RG_{\left[\psi(x),+\infty\right[}\Rder{\psi}_{\ast}F\right)_{\psi(x)}=0
  \]for each \(j\leqq k\). 
  Then we can apply Lemma \ref{Lem:1d-t-morse} and have
  \[
    H^{j}(\left]-\infty,b\right[;\Rder{\psi_{\ast}}F)
    \simarr
    H^{j}(\left]-\infty,a\right[;\Rder{\psi_{\ast}}F)
  \] for all \(j<k\).
  Then we can conclude that
  \[
    H^{j}(M_b\mathrel{;}F)
    \simarr
    H^{j}(M_a\mathrel{;}F)
  \]
  for all \(j<k\).
  \end{proof}

\section{Truncated Noncharacteristic Deformation}

We will reformulate the noncharacteristic deformation 
lemma \cite[Proposition 2.7.2]{KS90} to make a relation 
to truncated microsupport. 

\begin{Proposition+}\label{prop:t-nonchar}
  Let \(X\) be a Hausdorff space, \(F\in\Dompb(\zz_X)\), and 
  let \(\left(U_{t}\right)_{t\in\rr}\) be a family of 
  open subsets of \(X\). 
  We assume the following conditions. 
  \begin{enumerate}[(i)]
    \item \label{item:t-nonchar1}
    \(U_{t}=\bigcup_{s<t}U_{s}\), for all \(t\in\rr\). 
    \item \label{item:t-nonchar2}
    For all pairs \((s,t)\) with \(s\leqq t\), 
    the set \(\overline{U_{t}\smi U_{s}}\cap \supp{F}\) is compact.
    \item \label{item:t-nonchar3}
    Setting \(
      Z_{s}\coloneqq \bigcap_{t>s}\overline{U_t\smi U_s}
    \), we have for all pairs \((s,t)\) with \(s\leqq t\), 
    and all \(x\in Z_{s}\smi U_{t}\), \[
      (H^{j}_{X\smi U_{t}}(F))_{x}=0
      \quad 
      \text{for all \(j\leqq k\).}
    \]
    \item \label{item:t-nonchar4} 
    For any \(s\in\rr\), 
    \(\bigcap_{t>s}U_{t}=\overline{U_{s}}\) 
    and 
    \(\overline{U_{s}}\) is closed 
    in a paracompact open subset of \(X\).
    \item\footnote{
      This condition might be 
      deduced from \eqref{item:t-nonchar3}.
    } \label{item:t-nonchar5} 
    \(Z_s\smi U_{s}\subset\bd{U_{s}}\) 
    for all \(s\in\rr\). 
  \end{enumerate}
  Then, for all \(t\in\rr\), we have the isomorphisms\[
    H^{j}\left(\bigcup_{s}U_{s};F\right)
    \simar 
    H^{j}\left(U_{t};F\right)
    \quad \text{for all \(j< k\),} 
  \]
  and the injections\[
    H^{k}\left(\bigcup_{s}U_{s};F\right)
    \hookrightarrow 
    H^{k}\left(U_{t};F\right).
  \] 
\end{Proposition+}

\begin{proof}
  Consider the following conditions for 
  \(j\in\zz\), \(s\), and \(t\):
  \begin{align}
      \tag*{\((a)^j_s\)}\label{eq:t-nonchar-a}
      \mu_{s}^{j}&\colon
      \indlim[t>s] H^{j}(U_t;F) \simar H^{j}(U_s;F),\\
      \tag*{\((a')^k_s\)}\label{eq:t-nonchar-a'}
      \mu_{s}^{k}&\colon
      \indlim[t>s] H^{k}(U_t;F) \hookrightarrow H^{k}(U_s;F),\\
      \tag*{\((b)^j_s\)}\label{eq:t-nonchar-b}
      \lambda_{s}^{j}&\colon
      H^{j}(U_t;F)\simarr \prolim[s<t] H^{j}(U_s;F).
  \end{align}

  Let us prove \ref{eq:t-nonchar-a} for \(s\in \rr\) and \(j<k\). 
  We can assume that \(\overline{U_{t}\smi U_{s}}\) is compact 
  for any \(s\leqq t\) by replacing \(X\) by \(\supp{F}\). 
  Consider the distinguished triangle\[
    \RG_{X\smi U_{t}}(F)
    \to
    F
    \to
    \RG_{U_{t}}(F)
    \overset{+1}{\longrightarrow}.
  \]
  Applying to this triangle the functor \(
    (\blk)_{\overline{U_{t}}}
  \), we have \[
    \left(\RG_{X\smi U_{t}}(F)\right)_{\overline{U_{t}}}
    \to
    F_{\overline{U_{t}}}
    \to
    \RG_{U_{t}}(F)
    \overset{+1}{\longrightarrow}.
  \]
  Again applying to this triangle the functor \(\RG(X;\blk)\), 
  we have
  \begin{equation}
    \RG\left(X;\left(\RG_{X\smi U_{t}}(F)\right)_{\overline{U_{t}}}\right)
    \to
    \RG(\overline{U_{t}};F)
    \to
    \RG(U_{t};F)
    \overset{+1}{\longrightarrow}.
  \end{equation}
  By hypotheses \eqref{item:t-nonchar3}--\eqref{item:t-nonchar5}, 
  it follows that
  \[
    \left(H^{j}_{X\smi U_{t}}(F)\right)_{\bd{U_{t}}}=0
  \]
  for any \(j\leqq k\) and then \[
    H^{j}\left(X;
      \left(
        \RG_{X\smi U_{t}}(F)
      \right)_{\overline{U_{t}}}
    \right)=0
    \quad
    \text{for any \(j\leqq k\).}
  \] \footnote{この辺の議論はきちんと詰めないといけない．}
  Thus for any \(t\), we have
  \begin{align}
    \label{eq:a-iso}H^{j}(\overline{U_{t}};F)
    &\simar
    H^{j}(U_{t};F)\quad
    \text{for any \(j< k\),}\\
    \label{eq:a-inj}H^{k}(\overline{U_{t}};F)
    &\hookrightarrow
    H^{k}(U_{t};F).
  \end{align}
  Therefore, we have by \eqref{item:t-nonchar4} and \eqref{eq:a-iso},  
  \begin{align*}
    \indlim[t>s]H^{j}(U_t;F)
    &\cong 
    H^{j}\left(\bigcap_{t>s}U_t;F\right)\\
    &=
    H^{j}(\overline{U_s};F)\\
    &\simar
    H^{j}(U_s;F)
  \end{align*}
  for any \(j<k\), and \ref{eq:t-nonchar-a} follows. 
  Similarly, 
  we have by \eqref{item:t-nonchar4} and \eqref{eq:a-inj},
    \begin{align*}
    \indlim[t>s]H^{k}(U_t;F)
    &\cong 
    H^{k}\left(\bigcap_{t>s}U_t;F\right)\\
    &=
    H^{k}(\overline{U_s};F)\\
    &\hookrightarrow
    H^{k}(U_s;F)
  \end{align*}
  and \ref{eq:t-nonchar-a'} follows.
\end{proof}













\clearpage
\begin{Remark+}
  Lemma \ref{Lem:1d-t-morse} follows from 
  Proposition \ref{prop:t-nonchar}. 
  Indeed, putting \(X=\lopen-\infty,b\ropen\), 
  \(F=\mres[G]{I_{b}}\), and \(U_t=I_t\), 
  we find that \begin{enumerate}[(i)]
    \item \(\bigcup_{s<t} I_s
    = \bigcup_{s<t}\lopen-\infty,s\ropen
    = \lopen-\infty,t\ropen= I_t\) for all \(t\in\rr\). 
    \item If \(s\leqq t\), then \(
      \overline{I_t\smi I_s}\cap\supp{G}
      =\overline{\lclosed s,t\ropen}\cap\supp{G}
      =\lclosed s,t\rclosed\cap\supp{G}
      \) is compact.
    \item For \(s\in\rr\), we have \[
      Z_{s}
      =\bigcap_{t>s}\overline{I_{t}\smi I_{s}}
      =\bigcap_{t>s}\overline{\lclosed s,t\ropen}
      =\bigcap_{t>s}\lclosed s,t\rclosed
      =\left\{s\right\},
    \]
    and \[
      Z_s\smi I_t
      =\left\{s\right\}\smi \lopen-\infty,t\ropen
      =\begin{cases*}
        \left\{s\right\}& if \(s=t\),\\
        \varnothing& if \(s<t\).
      \end{cases*}
    \]
    By hypothesis of Lemma \ref{Lem:1d-t-morse}, 
    for all \(s\in \left\{s\right\}\), 
    \[
      H^{j}_{X\smi I_s}(F)_{s}
      =H^{j}_{\lclosed s,\infty\ropen}(F)_{s}=0.
    \]
    Then we can apply Proposition \ref{prop:t-nonchar}, and 
    we have the isomorphisms for all \(t\in \lclosed a,b\ropen\) and \(j<k\),
    \[
      H^{j}(I_{b};F)
      =H^{j}(\bigcup_{s}I_{s};F)
      \simar
      H^{j}(I_{t};F)
    \]
    and the injections
    \[
      H^{k}(I_{b};F)
      =H^{k}(\bigcup_{s}I_{s};F)
      \hookrightarrow
      H^{k}(I_{t};F)
    \]
  \end{enumerate}
\end{Remark+}


\section{
  Relative Microlocal Morse Theory via Truncated Microsupport
}

Let \(f\colon Y\to X\) be a morphism of manifolds, and 
\(\varphi\) a real valued function of class \(C^1\). 
For a real number \(t\in \rr\), we set 
\begin{align*}
  Y_t
  &\coloneqq \left\{y\in Y\smid \varphi(y)< t\right\} 
  = \varphi^{-1}(\lopen-\infty,t\ropen),\\
  \widetilde{Y_t}
  &\coloneqq \left\{y\in Y\smid \varphi(y)\leqq t\right\} 
  = \varphi^{-1}(\lopen-\infty,t\rclosed).
\end{align*}
Let us define the following morphisms:
\[\xymatrix{
  Y_t\ar[r]^-{j_t}
  \ar[rd]_-{f_t}
  &
  Y
  \ar[d]^-{f}\\
  &X,
}\quad\xymatrix{
  \widetilde{Y_t}\ar[r]^-{\tilde{j}_t}
  \ar[rd]_-{f_t}
  &
  Y
  \ar[d]^-{f}\\
  &X.
}\]

\begin{Proposition+}
  Let \(G\in\Dompb(Y)\), and assume \(
    f_{t}\colon \supp(G)\cap \overline{Y_t}\to X
  \) is proper for all \(t\in\rr\). 
  Let \(k\) be an integer and \(t_\circ\) a real number.
  \begin{enumerate}[(i)]
    \item Assume that, 
    for any \(y\in Y\smi \widetilde{Y}_{t_{\circ}}\), 
    we have \[
      d\varphi(y)\notin \left(
        \muS[k](G)+ \mdiff[f]\left(
          Y\pb[X]\cotb[X]
        \right)
      \right).
    \]
    Then \begin{align}
      \tag{a}
      \Rder[j]{f_{\ast}}G
      &\cong
      \Rder[j]{{f_{t}}_{\ast}}j^{-1}_{t}G
      \quad
      \text{ for \(j<k\) and \(t>t_\circ\)},\\
      \tag{a\(_k\)}
      \Rder[k]{f_{\ast}}G
      &\hookrightarrow
      \Rder[k]{{f_{t}}_{\ast}}j^{-1}_{t}G
      \quad
      \text{ for \(t>t_\circ\)},\\
      \tag{a\('\)}
      \Rder[j]{\tilde{f}_{\ast}}G
      &\cong
      \Rder[j]\tilde{f_{t}}_{\ast}\tilde{j}^{-1}_{t}G
      \quad
      \text{ for \(j<k\) and \(t\geqq t_\circ\)},\\
      \tag{a\(_k'\)}
      \Rder[k]{\tilde{f}_{\ast}}G
      &\hookrightarrow
      \Rder[k]\tilde{f_{t}}_{\ast}\tilde{j}^{-1}_{t}G
      \quad
      \text{ for \(t\geqq t_\circ\)},\\
      \tag{b}
      \muS[k](\Rder[]{f_{\ast}}G)
      &\subset
      \mpb[f]\left(\mdiff[f]^{-1}\left(
        \muS[k](G)\cap \pi^{-1}(\widetilde{Y}_{t_{\circ}})
      \right)\right)\\
      %&
      \notag&=\mpb[f]\mdiff[f]^{-1}\left(
        \muS[k](G)
      \right).
    \end{align}
    \item Assume that, 
    for any \(y\in Y\smi \widetilde{Y}_{t_{\circ}}\), 
    we have \[
      -d\varphi(y)\notin \left(
        \muS[k](G)+ \mdiff[f]\left(
          Y\pb[X]\cotb[X]
        \right)
      \right).
    \]
    Then \begin{align}
      \tag{c}
      \Rder[j]{f_{!}}G
      &\cong
      \Rder[j]{{f_{t}}_{!}}j^{-1}_{t}G
      \quad
      \text{ for \(j<k\) and \(t>t_\circ\)},\\
      \tag{c\(_k\)}
      \Rder[k]{f_{!}}G
      &\hookrightarrow
      \Rder[k]{{f_{t}}_{!}}j^{-1}_{t}G
      \quad
      \text{ for \(t>t_\circ\)},\\
      \tag{c\('\)}
      \Rder[j]{\tilde{f}_{!}}G
      &\cong
      \Rder[j]\tilde{f_{t}}_{\ast}\tilde{j}^{!}_{t}G
      \quad
      \text{ for \(j<k\) and \(t\geqq t_\circ\)},\\
      \tag{c\(_k'\)}
      \Rder[k]{\tilde{f}_{!}}G
      &\hookrightarrow
      \Rder[k]\tilde{f_{t}}_{\ast}\tilde{j}^{!}_{t}G
      \quad
      \text{ for \(t\geqq t_\circ\)},\\
      \tag{d}
      \muS[k](\Rder[]{f_{!}}G)
      &\subset
      \mpb[f]\left(\mdiff[f]^{-1}\left(
        \muS[k](G)\cap \pi^{-1}(\widetilde{Y}_{t_{\circ}})
      \right)\right)\\
      %&
      \notag&=\mpb[f]\mdiff[f]^{-1}\left(
        \muS[k](G)
      \right).
    \end{align}
  \end{enumerate}
\end{Proposition+}

\begin{proof}
  Let \(\tilde{f}\) be the morphism \((f,\varphi)\colon Y\to X\times \rr\), and let \(q\) and \(\tilde{\varphi}\) be the first and second projections defined on \(X\times \rr\) respectively.
\end{proof}
%===============================================
% 参考文献スペース
%===============================================
\begin{thebibliography}{20} 
  \bibitem[GKS12]{GKS12} St\'ephane Guillermou, 
  Masaki Kashiwara, Pierre Schapira, 
  \textit{Sheaf Quantization of Hamiltonian Isotopies and Applications to Nondisplaceability Problems}, 
  Duke. Math. J. 161, No. 2, pp.201--245, 2012.  
  \bibitem[KMS03]{KMS03} Masaki Kashiwara, Teresa Monteiro Fernandes, Pierre Schapira, 
    \textit{Truncated Microsupprt and Holomorphic Solutions 
    of \(D\)-modules}, 
    Ann. Scient. \'Ec. Norm. Sup., s\'erie 4, tome 36, 
    pp.583--399, 2003.
  \bibitem[KS90]{KS90} Masaki Kashiwara, Pierre Schapira, 
    \textit{Sheaves on Manifolds}, 
    Grundlehren der Mathematischen Wissenschaften, 292, Springer, 1990.
  \bibitem[ST92]{ST92} Pierre Schapira, Nobuyuki Tose, 
    \textit{Morse Inequalities for \(\rr\)-constructible Sheaves}, 
    Adv. Math. 93, pp.1--8, 1992. 
\end{thebibliography}

%===============================================




\end{document}
